\documentclass[11pt,a4paper]{article}
\usepackage[top=1.7cm,bottom=1.5cm,left=1.5cm,right=1.5cm,headsep=8pt,footskip=16pt]{geometry}
\usepackage{enumitem}
\usepackage{amssymb}
\usepackage{array}
\usepackage[table]{xcolor}
\usepackage{fancyhdr}
\usepackage{lastpage}
\usepackage{multicol}
\usepackage{polyglossia}
\setdefaultlanguage[calendar=gregorian, numerals=maghrib]{arabic}
\setotherlanguage{english}
\IfFontExistsTF{Amiri}{%
  \newfontfamily\arabicfont[Script=Arabic, Scale=1.08]{Amiri}
}{%
  \newfontfamily\arabicfont[Script=Arabic, Scale=1.08, Path=../fonts/, Extension=.ttf,
    UprightFont=*-Regular, BoldFont=*-Bold, ItalicFont=*-Italic,
    BoldItalicFont=*-BoldItalic]{Amiri}
}
\IfFontExistsTF{AlBattar}{%
  \newfontfamily\khatt[Script=Arabic, Scale=1.15]{AlBattar}
}{%
  \newfontfamily\khatt[Script=Arabic, Scale=1.15, Path=../fonts/,
    BoldFont=ae_AlBattar.ttf]{ae_AlBattar.ttf}
}
\renewcommand{\headrulewidth}{0.8pt}
\renewcommand{\footrulewidth}{0.8pt}
\fancyhf{}
\rhead{\khatt المدرسة العليا للأساتذة ببوزريعة}
\lhead{\khatt ملحقة بومرداس}
\lfoot{\khatt الأستاذ: د. بن عبد الرزاق يوسف}
\cfoot{\thepage/\pageref{LastPage}}
\pagestyle{fancy}
\parindent 0pt
\setlength{\columnsep}{0.6cm}
\setlength{\columnseprule}{0.4pt}
\setlist[enumerate,1]{label=\textbf{\arabic*.}, leftmargin=*, itemsep=5pt, topsep=2pt}
\begin{document}

\linespread{1.08}\selectfont
\begin{center}
{\LARGE\khatt\textbf{الاختبار الثالث عشر في الذكاء الاصطناعي}}\\[4pt]
{\normalsize وضعيات إشكالية ــ أسئلة متعددة الاختيارات ــ الفصول من 1 إلى 6}
\end{center}
\vspace{-2pt}\rule{\textwidth}{0.6pt}\\[2pt]
\textbf{التعليمة:} يتضمن الاختبار خمس وضعيات إشكالية. اقرأ كل وضعية بعناية، ثم اختر التحليل أو التصرف الأنسب
من بين الاختيارات الخمسة؛ لكل وضعية إجابة صحيحة واحدة فقط.
اكتب \textbf{حرف} الإجابة الصحيحة في جدول الإجابات في آخر الاختبار.
لكل وضعية 4 نقاط (المجموع 20 نقطة).
\vspace{4pt}

\noindent
{\large\khatt\textbf{الوضعية 1}}\hfill{\small(الفصل 1)}\par\nopagebreak\vspace{3pt}
\setlength{\fboxsep}{7pt}\colorbox{gray!10}{\parbox{\dimexpr\linewidth-14pt}{تستعمل أستاذة أداة توليدية لإنتاج تمارين وتصحيحها، فلاحظت بعد أسابيع أن مفتاح التصحيح المرفق يتضمن أحيانًا إجابة خاطئة لتمرين صحيح الصياغة، مع أن الأداة «هي التي وضعت التمرين». ويرى زميل أن ذلك مستحيل. ويريد الزميل أن يكتفي بإعادة الطلب كلما شك في المفتاح. ما التفسير والتصرف الأنسب؟}}\par\nopagebreak
\begin{itemize}[leftmargin=1.0cm, labelsep=6pt, itemsep=3pt, topsep=4pt, parsep=0pt]
  \item[\textbf{أ)}] الخطأ سببه أن المفتاح طُلب في المحادثة نفسها، ولو طُلب في محادثة جديدة لجاء صحيحًا بالضرورة، فتكفي إعادة الطلب في محادثة جديدة دون مراجعة المفتاح.
  \item[\textbf{ب)}] الأداة تولّد التمرين والمفتاح توليدًا احتماليًا، والحل أن تطلب منها في نهاية كل مرة «تأكد من صحة المفتاح»، ثم تعتمده دون مراجعة، لأن هذه العبارة تضمن الاتساق.
  \item[\textbf{ج)}] زميلها محق، فالأداة لا تخطئ في مفتاح تمرين وضعته بنفسها، والخطأ الظاهر سببه أن الأستاذة فهمت التمرين خطأ، فتعتمد المفتاح كما هو في كل مرة.
  \item[\textbf{د)}] الأداة تولّد التمرين والمفتاح كلًّا على حدة توليدًا احتماليًا، ولا تضمن الاتساق بينهما؛ فتحلّ الأستاذة التمارين بنفسها أو تراجع المفتاح سؤالًا سؤالًا قبل استعماله.
  \item[\textbf{هـ)}] الأداة تخطئ في المفتاح عمدًا لتدفع المستخدم إلى المراجعة، لأنها مصممة تربويًا لتنمية اليقظة، ولذلك تُبقي الأستاذة الخطأ في المفتاح لتختبر به التلاميذ.
\end{itemize}
\par\vspace{6pt}
\noindent
{\large\khatt\textbf{الوضعية 2}}\hfill{\small(الفصل 3)}\par\nopagebreak\vspace{3pt}
\setlength{\fboxsep}{7pt}\colorbox{gray!10}{\parbox{\dimexpr\linewidth-14pt}{تريد أستاذة نشاطًا حواريًا يحاور فيه التلاميذ «ابن خلدون» عبر أداة ذكاء اصطناعي تتقمص شخصيته، لتقريب مقدمته إليهم. وتخشى أن تنسب الأداة إليه أقوالًا لم يقلها، أو أن تجيب عن أسئلة معاصرة كأنه يعرفها. والحصة مدتها ساعة واحدة، والتلاميذ متحمسون للفكرة. أي إعداد للطلب أنسب لهذا النشاط؟}}\par\nopagebreak
\begin{itemize}[leftmargin=1.0cm, labelsep=6pt, itemsep=3pt, topsep=4pt, parsep=0pt]
  \item[\textbf{أ)}] تكتب «أنت ابن خلدون» فقط، لأن الأداة تعرف كل أقواله، وتقمص الشخصية يجعلها تلتزم تلقائيًا بأفكاره وعصره، دون حاجة إلى قيود أو تجربة مسبقة.
  \item[\textbf{ب)}] تحدد الدور والجمهور، وتطلب الاقتصار على أفكار المقدمة المعروفة مع التنبيه إلى أن الحوار تمثيل تعليمي، وأن ترد الأداة على الأسئلة المعاصرة بأنها خارج عصره، ثم تجرّب الحوار وتراجعه قبل الحصة.
  \item[\textbf{ج)}] تتجنب تحديد الجمهور، لأن ذكر أن المحاورين تلاميذ يجعل الأداة تبسّط أفكار ابن خلدون حتى تفقد دقتها، والأفضل أن تخاطبهم بلغة المقدمة الأصلية.
  \item[\textbf{د)}] تحدد الدور والقيود جيدًا، ثم تترك التلاميذ يحاورون الأداة دون تجربة مسبقة، ويعدّون كل ما تقوله أقوالًا موثقة لابن خلدون يستشهدون بها في بحوثهم، فيكتسبون ثقة في أنفسهم وفي أدوات البحث الحديثة.
  \item[\textbf{هـ)}] تطلب من الأداة أن تجيب كابن خلدون عن كل سؤال معاصر بثقة، لأن ذلك يجعل الحوار أكثر تشويقًا، ويُقنع التلاميذ بأن أفكاره صالحة لكل المسائل الحديثة، ويرسّخ في أذهانهم قيمة التراث وصلاحيته.
\end{itemize}
\par\vspace{6pt}
\noindent
{\large\khatt\textbf{الوضعية 3}}\hfill{\small(الفصل 4)}\par\nopagebreak\vspace{3pt}
\setlength{\fboxsep}{7pt}\colorbox{gray!10}{\parbox{\dimexpr\linewidth-14pt}{نشر أستاذ على صفحة المؤسسة صورًا مولَّدة لمعركة تاريخية جزائرية مشهورة، تبدو كأنها صور فوتوغرافية حقيقية، دون أي إشارة إلى أنها مولَّدة. فتداولها كثيرون على أنها وثائق أصلية من الأرشيف. ويطالب بعض المتابعين بحذف المنشور دون توضيح، ويرى آخرون الإبقاء عليه لأن هدفه نبيل. ما التقدير الأخلاقي والتصرف الأنسب؟}}\par\nopagebreak
\begin{itemize}[leftmargin=1.0cm, labelsep=6pt, itemsep=3pt, topsep=4pt, parsep=0pt]
  \item[\textbf{أ)}] يحذف الصور بصمت دون أي توضيح، لأن الاعتراف بالخطأ يسيء إلى سمعة المؤسسة، والحذف وحده يكفي لإيقاف تداولها على أنها وثائق أرشيفية أصلية.
  \item[\textbf{ب)}] لا إشكال، لأن الصور تخدم الذاكرة الوطنية، والجمهور يعرف أن التصوير الفوتوغرافي للمعارك القديمة نادر، فلا يحتاج الأستاذ إلى أي تنبيه أو تصحيح.
  \item[\textbf{ج)}] المسؤولية على من تداولها، لا على الأستاذ، لأنه لم يقل إنها حقيقية، والسكوت عن طبيعة الصورة لا يُعدّ تضليلًا ما لم يصرّح الناشر بمعلومة خاطئة صراحة.
  \item[\textbf{د)}] يُبقي الصور ويضيف تعليقًا يؤكد أنها من الأرشيف، لأن الصور المولَّدة تمثّل الحدث بدقة، والمهم أن يصل المعنى التاريخي إلى الجمهور بغض النظر عن مصدرها.
  \item[\textbf{هـ)}] نشر صور مولَّدة واقعية دون تنبيه يضلل الجمهور ويخلط التاريخ بالمتخيَّل؛ فيضيف تنبيهًا واضحًا بأنها مولَّدة، ويصحّح المعلومة علنًا، ويلتزم مستقبلًا بالتصريح عند كل نشر.
\end{itemize}
\par\vspace{6pt}
\noindent
{\large\khatt\textbf{الوضعية 4}}\hfill{\small(الفصل 5)}\par\nopagebreak\vspace{3pt}
\setlength{\fboxsep}{7pt}\colorbox{gray!10}{\parbox{\dimexpr\linewidth-14pt}{وجد طالب على موقع غير معروف «إحصائية» تقول إن 70\% من الجزائريين لا يقرؤون كتابًا في السنة، ثم سأل أداة ذكاء اصطناعي فأكدتها وأضافت أنها «من دراسة لليونسكو». ولم يجد الدراسة بالبحث في موقع اليونسكو. ويريد أن يستشهد بها في مقدمة مذكرته عن المطالعة. كيف يتعامل مع هذه الإحصائية؟}}\par\nopagebreak
\begin{itemize}[leftmargin=1.0cm, labelsep=6pt, itemsep=3pt, topsep=4pt, parsep=0pt]
  \item[\textbf{أ)}] يعدّل الرقم إلى 60\% ليكون أكثر اعتدالًا ثم يعتمده، لأن الإحصاءات المتداولة مبالغ فيها عادة، والتعديل الحذر يجعلها أقرب إلى الواقع دون حاجة إلى مصدر.
  \item[\textbf{ب)}] يبحث عن الرقم في محرك البحث، فإن وجده في مواقع كثيرة اعتمده ونسبه إلى اليونسكو، لأن كثرة تداوله دليل على أنه مأخوذ من دراسة حقيقية للمنظمة، ثم يذكر في الهامش عدد المواقع التي نشرته.
  \item[\textbf{ج)}] لا يعتمدها، لأن تأكيد الأداة قد يكون إعادة إنتاج لمعلومة متداولة لا تحققًا منها، ونسبتها إلى اليونسكو بلا إحالة دقيقة لا تكفي؛ فيبحث عن الدراسة الأصلية أو عن إحصاءات رسمية موثقة.
  \item[\textbf{د)}] يعتمدها مع ذكر الأداة مصدرًا، لأن الأداة أضافت معلومة المصدر، فتصبح الإحالة إليها إحالة غير مباشرة إلى اليونسكو، وهي مقبولة في البحوث الجامعية.
  \item[\textbf{هـ)}] يعتمدها، لأن اتفاق الموقع والأداة على الرقم نفسه دليل كافٍ على صحته، وذكر اليونسكو يمنحها مصداقية، حتى لو لم يجد الدراسة الأصلية في موقع المنظمة، ويكفي ذلك لتوثيقها في مقدمة المذكرة.
\end{itemize}
\par\vspace{6pt}
\noindent
{\large\khatt\textbf{الوضعية 5}}\hfill{\small(الفصل 6)}\par\nopagebreak\vspace{3pt}
\setlength{\fboxsep}{7pt}\colorbox{gray!10}{\parbox{\dimexpr\linewidth-14pt}{صوّرت أستاذة السبورة في نهاية حصة نحو، وطلبت من أداة متعددة الأنماط تحويل الصورة إلى ملخص منظم للدرس ترسله إلى التلاميذ الغائبين. جاء الملخص جيدًا، لكن الأداة أضافت أمثلة لم تُكتب على السبورة، وأخطأت في قراءة سهم يربط قاعدة بمثالها. وبعض التلاميذ الغائبين سيعتمدون على الملخص وحده في المراجعة. ما التصرف الأنسب؟}}\par\nopagebreak
\begin{itemize}[leftmargin=1.0cm, labelsep=6pt, itemsep=3pt, topsep=4pt, parsep=0pt]
  \item[\textbf{أ)}] تقابل الملخص بالسبورة وبما شرحته فعلًا، فتصحح الربط الخاطئ، وتراجع الأمثلة المضافة فتحذفها أو تُبقي الصحيح منها مع التأكد من مناسبته، لأن الملخص سيصبح مرجع الغائبين.
  \item[\textbf{ب)}] ترسل صورة السبورة وحدها دون ملخص، لأن الأدوات متعددة الأنماط لا تفهم ما يُكتب بخط اليد إطلاقًا، وكل ما ورد في الملخص مولَّد لا علاقة له بالصورة.
  \item[\textbf{ج)}] تحذف كل الأمثلة من الملخص، بما فيها أمثلة السبورة، لأن وجود مثال مضاف واحد يجعل كل الأمثلة موضع شك، ثم ترسل القواعد وحدها دون أي تصحيح للربط.
  \item[\textbf{د)}] ترسل الملخص كما هو، لأن الأمثلة المضافة تثري الدرس، والأداة تقرأ الأسهم والروابط بدقة تامة، فما تراه الأستاذة خطأ في الربط تفسير مختلف مقبول للسبورة.
  \item[\textbf{هـ)}] تطلب من الأداة إعادة الملخص، وترسل النسخة الثانية دون مقابلة، لأن الأداة تتنبّه في المحاولة الثانية إلى أخطاء القراءة الأولى، فتصحح الربط وتحذف الإضافات تلقائيًا.
\end{itemize}
\par\vspace{6pt}
\begin{center}
{\khatt\textbf{جدول الإجابات}}\\[3pt]
\renewcommand{\arraystretch}{1.6}
\begin{tabular}{|>{\centering\arraybackslash}m{1.9cm}|*{5}{>{\centering\arraybackslash}m{1.2cm}|}}
\hline
\textbf{رقم السؤال} & \textbf{1} & \textbf{2} & \textbf{3} & \textbf{4} & \textbf{5} \\ \hline
\textbf{الإجابة} &  &  &  &  &  \\[8pt] \hline
\end{tabular}

\vspace{8pt}
\renewcommand{\arraystretch}{2.2}
\begin{tabular}{|>{\centering\arraybackslash}m{7.5cm}|>{\centering\arraybackslash}m{3.2cm}|>{\centering\arraybackslash}m{3.2cm}|}
\hline
\textbf{الاسم واللقب} & \textbf{الفوج} & \textbf{العلامة} \\ \hline
 & & \hfill /20 \\ \hline
\end{tabular}
\end{center}
\end{document}
